\chapter{Introdução}

Em dezembro de 1906, o governo brasileiro criou uma instituição autorizada a receber ouro e moedas estrangeiras conversíveis e a entregar, em troca, notas resgatáveis à taxa de 15 pence por mil-réis. A Caixa de Conversão limitou a valorização da moeda, ampliou a circulação conforme o ouro ingressava e aproximou o Brasil das regras do padrão-ouro sem converter imediatamente todo o papel-moeda existente. Em dezembro de 1910, a taxa foi elevada a 16 pence. Em agosto de 1914, diante da ruptura dos mercados provocada pela guerra, o troco das notas foi suspenso. Entre esses marcos, a Caixa ocupou lugar recorrente na política econômica, no noticiário financeiro e nas controvérsias sobre câmbio, café, crédito e relações com banqueiros estrangeiros.

A instituição nasceu num ponto de inflexão. A política de contração, equilíbrio fiscal e valorização cambial adotada depois do \textit{funding loan} de 1898 havia restaurado a confiança externa e reduzido parte das pressões que marcaram a primeira década republicana. Ao mesmo tempo, a valorização do mil-réis diminuiu as receitas domésticas dos exportadores, num momento de excesso de oferta e queda dos preços do café. O Convênio de Taubaté reuniu medidas de defesa do produto e uma proposta de estabilização cambial. Quando os projetos foram separados, a Caixa deixou de ser apenas cláusula de um acordo cafeeiro e passou a concentrar uma discussão nacional sobre o nível e o desenho da estabilidade.

Esta dissertação investiga como quatro jornais organizaram esse debate entre 1906 e 1914: \textit{O Paiz}, \textit{Correio da Manhã}, \textit{Correio Paulistano} e \textit{Gazeta de Notícias}. Pergunta-se quais objetos de política foram discutidos, quais atores e grupos receberam voz, que categorias e termos foram mobilizados e como a experiência do regime alterou o conteúdo do debate. O problema não se limita a classificar cada periódico como favorável ou contrário à Caixa. Uma mesma instituição podia ser defendida ou criticada por sua taxa, pelo limite de emissão, pela composição do lastro, pela relação com o Banco do Brasil, pela competência do Congresso, pelo destino dos fundos monetários ou pelos custos impostos a determinados interesses.

O recorte começa em 1906, quando o Convênio, os projetos e a lei tornam a Caixa objeto definido, e termina em 1914, quando a conversibilidade é suspensa sob uma crise internacional de natureza extraordinária. A cronologia permite observar quatro situações distintas: criação, operação inicial, reforma e crise. As categorias que discriminavam posições em 1906 não necessariamente conservaram o mesmo sentido em 1910 ou 1914. A comparação temporal é, portanto, parte do argumento, e não apenas extensão do número de fontes.

Os jornais formam o corpus principal porque constituíram espaços de circulação, seleção e elaboração do debate. O estudo os complementa com debates parlamentares, legislação, mensagens de governo, cartas, telegramas, relatórios financeiros e as edições anuais do \textit{Retrospecto Commercial}. Essas fontes acessórias ajudam a reconstruir decisões e a distinguir documento reproduzido de comentário editorial. Elas enriquecem a interpretação, mas não deslocam a pergunta central sobre o que os periódicos publicaram, destacaram e assumiram como voz própria.

\interpretacaoprovisoria{A hipótese central é que, em 1906, estabilidade e conversibilidade já possuíam ampla legitimidade, enquanto o conflito principal incidia sobre o nível da taxa, o ritmo da valorização, a emissão, o lastro, as competências institucionais e os custos distributivos. O debate não se organizava simplesmente como continuação da clivagem entre metalismo e papelismo. Entre 1912 e 1914, a reversão externa separou três questões antes mais facilmente sobrepostas: se a estabilidade continuava desejável, se a conversibilidade permanecia viável e se o custo de defender a taxa era aceitável.}

Essa hipótese não supõe consenso absoluto. Em 1906 havia defensores do retorno a 27 pence, partidários da fixação a 15 e proponentes de um novo padrão em torno de 12. O ponto é que mesmo propostas de taxa baixa podiam invocar moeda metálica, conversibilidade e saneamento. A oposição formal à Caixa também reunia razões incompatíveis. Por isso, categorias como ortodoxo, expansionista e metalista serão tratadas como construções históricas e historiográficas que precisam ser decompostas em posições observáveis.

A contribuição documental é inventarial. O projeto preserva 11.960 objetos digitais dos quatro jornais, com 117.703 páginas dotadas de texto. A menção tolerante a ruído localizou 8.331 páginas nas quais a Caixa aparece, mas grande parte corresponde a cotações, boletins e rotina. O desafio é descrever, da maneira mais exaustiva possível dentro das restrições do acervo e das ferramentas, onde o debate substantivo ocorre, quem participa dele e como suas categorias mudam. A quantificação será usada para caracterizar cobertura e composição, enquanto o argumento dependerá da leitura contextualizada das peças e da conferência das passagens decisivas na imagem.

A contribuição interpretativa consiste em aproximar duas literaturas que nem sempre são mobilizadas em conjunto. A história econômica explica a Caixa como resposta à trajetória monetária, à crise do café e à inserção periférica no padrão-ouro. A história da imprensa exige examinar o periódico como instituição, empresa e espaço de vozes. Ao confrontar as categorias da historiografia com o vocabulário e os gêneros das fontes, a dissertação procura identificar tanto correspondências quanto inadequações. Um discurso parlamentar transcrito, por exemplo, documenta uma posição presente nas páginas do jornal, mas não prova endosso editorial. Um editorial contrário à elevação da taxa pode defender a Caixa existente e, ao mesmo tempo, resistir a nova valorização.

O texto está organizado em cinco capítulos, além desta introdução e da conclusão. O Capítulo 1 apresenta a imprensa como fonte, delimita o corpus e descreve o procedimento de inventário, catalogação e leitura. O Capítulo 2 parte do padrão-ouro internacional e reconstrói os antecedentes brasileiros, da crise monetária da primeira década republicana ao Convênio de Taubaté. O Capítulo 3 acompanha a criação da Caixa em 1906 e decompõe as controvérsias sobre taxa, emissão, fundos, autoridade e agência externa. O Capítulo 4 examina os primeiros anos de funcionamento, a crise de 1907, a recuperação das entradas de ouro e a reforma de 1910. O Capítulo 5 analisa a custódia dos depósitos, a reversão iniciada antes da guerra e a suspensão de 1914. A conclusão retorna à trajetória de cada periódico e explicita o que a documentação permite rever nas categorias da literatura.

O resultado aqui apresentado é um rascunho de trabalho. Passagens marcadas como fontes a conferir dependem de inspeção final na imagem do periódico; interpretações assinaladas como provisórias indicam hipóteses que ainda precisam sobreviver à leitura comparada de todo o material substantivo. Tornar essas pendências visíveis é parte do método e evita que a fluidez da redação esconda graus diferentes de certeza documental.

